\begin{codebox}
\codeline{\textcolor{cmtgray}{\#\ ==============================}}
\codeline{\textcolor{cmtgray}{\#\ 5.\ 工程化使用建议}}
\codeline{\textcolor{cmtgray}{\#\ ==============================}}
\codeline{\textcolor{cmtgray}{\#\ 不必对全本体做完整标注——成本高；}}
\codeline{\textcolor{cmtgray}{\#\ 推荐在以下时机做局部OntoClean审查：}}
\codeline{\textcolor{cmtgray}{\#\ \ \ (1)\ 类层次初版评审时：只标注\ R（刚性），揪出"角色当类型"}}
\codeline{\textcolor{cmtgray}{\#\ \ \ (2)\ 两个本体合并时：检查同一性判据是否冲突}}
\codeline{\textcolor{cmtgray}{\#\ \ \ (3)\ 推理出现反直觉结果时：用元属性定位概念混淆}}
\codeline{\textcolor{cmtgray}{\#}}
\codeline{\textcolor{cmtgray}{\#\ 经验法则（覆盖80\%问题）：}}
\codeline{\textcolor{cmtgray}{\#\ \ \ 会"变回去"的概念（状态/角色/阶段）不要放进\ rdfs:subClassOf\ 层次，}}
\codeline{\textcolor{cmtgray}{\#\ \ \ 用属性（hasStatus\ /\ playsRole\ /\ hasPhase）表达}}
\end{codebox}
